\section{成化至正德的补充编年节点（二）}
以下节点继续沿同一编年脉络展开。每一锚点均紧接可供核查的事件叙述；来源能确认何种行动、文书或变动，不能自动消除不同记录之间的立场差异。

\subsection{1506年前后的续接}

\eventanchor{ML-1506-004}{明·1506（正德元年）}
\paragraph{1506（正德元年）：正德皇帝乔装游走北京街头} 这项记录把本章所见的制度运作落实到具体年份。其形成背景是：继承、官僚协调或皇权运作的压力使问题进入朝廷程序。 它带来的、而且仍可由后续材料追踪的改变是：它影响后续人事与政策议程；动机和派系标签不能由单一叙事裁定。 《武宗实录》正德元年条；《明史·本纪》及相关列传；诏令、奏疏或会典版本 提供相互检验的起点；但译写候选以编年材料定位；实录记载反映朝廷选择，崇祯期须另以奏疏、地方及敌对政权材料交叉。 因此本段仅据材料确认事件的时序、行动和可见后果，并把政策文本与地方实践、报称数字和社会经历分开处理。

\eventanchor{MM-1506-EIGHT-TIGERS}{明·1506（正德元年）}
\paragraph{1506（正德元年）：刘瑾等八名宦官获武宗信任并进入中枢政治} 这项记录把本章所见的制度运作落实到具体年份。其形成背景是：新君年少，内廷近侍在日常起居和诏令传达中占据接近皇权的位置。 它带来的、而且仍可由后续材料追踪的改变是：刘瑾集团得以影响人事、财赋和厂卫案件；“八虎”这一标签本身也塑造了后世的道德化叙事。 《武宗实录》正德元年条；《明史·刘瑾传》；《明史纪事本末·刘瑾用事》 提供相互检验的起点；但“八虎”称谓及成员主要见于后出叙事；近侍受宠和刘瑾权势上升可由实录、传记互证。 因此本段仅据材料确认事件的时序、行动和可见后果，并把政策文本与地方实践、报称数字和社会经历分开处理。

\eventanchor{ML-1507-001}{明·1507（正德二年）九月}
\paragraph{1507（正德二年）九月：元宵节花灯费35万两银子} 这项记录把本章所见的制度运作落实到具体年份。其形成背景是：财政征解、交通工程或货币支付的压力使相关措施进入政务。 它带来的、而且仍可由后续材料追踪的改变是：其法定安排不等于地方执行，后续须以地域材料追踪负担与成效。 《武宗实录》正德二年条；《明会典》相关条目；地方志、黄册/契约/河工材料 提供相互检验的起点；但译写候选以编年材料定位；实录记载反映朝廷选择，崇祯期须另以奏疏、地方及敌对政权材料交叉。 因此本段仅据材料确认事件的时序、行动和可见后果，并把政策文本与地方实践、报称数字和社会经历分开处理。

\eventanchor{ML-1507-002}{明·1507（正德二年）}
\paragraph{1507（正德二年）：科举、讲学和官员文集的本年线索可用于追踪知识网络，主张与行政结果须分开。} 这项记录把本章所见的制度运作落实到具体年份。其形成背景是：制度、教育或知识传播的需要推动了相关行动或文本形成。 它带来的、而且仍可由后续材料追踪的改变是：它提供制度和传播史的锚点，不能据此推断社会整体接受。 《武宗实录》正德二年条；《明史·选举志》或《艺文志》；文集、碑刻或书目材料 提供相互检验的起点；但桥接条目只标示可回查的年度问题与材料链；实录有中央选择性，崇祯期尤其需要奏疏、地方、朝鲜及满文材料交叉。 因此本段仅据材料确认事件的时序、行动和可见后果，并把政策文本与地方实践、报称数字和社会经历分开处理。

\eventanchor{ML-1507-003}{明·1507（正德二年）}
\paragraph{1507（正德二年）：嘉靖初继统礼制的本年文书构成大礼议过程锚点，礼制理由和政治效应不可简化为单一动机。} 这项记录把本章所见的制度运作落实到具体年份。其形成背景是：继承、官僚协调或皇权运作的压力使问题进入朝廷程序。 它带来的、而且仍可由后续材料追踪的改变是：它影响后续人事与政策议程；动机和派系标签不能由单一叙事裁定。 《武宗实录》正德二年条；《明史·本纪》及相关列传；诏令、奏疏或会典版本 提供相互检验的起点；但桥接条目只标示可回查的年度问题与材料链；实录有中央选择性，崇祯期尤其需要奏疏、地方、朝鲜及满文材料交叉。 因此本段仅据材料确认事件的时序、行动和可见后果，并把政策文本与地方实践、报称数字和社会经历分开处理。

\eventanchor{MM-1507-LIU-JIN-FISCAL-EXTRACTIONS}{明·1507（正德二年）}
\paragraph{1507（正德二年）：刘瑾集团要求地方严核钱粮、差役与库藏，内廷对财政文书的介入加深} 这项记录把本章所见的制度运作落实到具体年份。其形成背景是：内廷供用、赏赐和军费增加，中央认为既有赋役账册存在隐漏，刘瑾集团借近侍渠道扩大核查权。 它带来的、而且仍可由后续材料追踪的改变是：地方官与胥吏须应对新的报送和催征压力；财政整饬也成为1510年清算刘瑾时指认苛扰的材料来源。 《武宗实录》正德二年户部与内官条；《明史·刘瑾传》；《明史·食货志》 提供相互检验的起点；但诏令、催核和中官参与可由实录及制度志确认；追征数额和实际负担受申报、归罪和账册口径影响，不能拼作全国实收。 因此本段仅据材料确认事件的时序、行动和可见后果，并把政策文本与地方实践、报称数字和社会经历分开处理。

\eventanchor{MM-1507-LIU-JIN-PURGE}{明·1507（正德二年）}
\paragraph{1507（正德二年）：刘瑾借厂卫和考察处置异己官员，朝廷人事与言路受到冲击} 这项记录把本章所见的制度运作落实到具体年份。其形成背景是：武宗即位后近侍取得接近皇权的优势，刘瑾需排除批评其财政和人事政策的官员以稳固集团。 它带来的、而且仍可由后续材料追踪的改变是：中枢批评渠道收缩，地方官迁调与案件相互牵连；其后刘瑾失势时，被压制的官僚网络成为反向清算力量。 《武宗实录》正德二年吏部与锦衣卫条；《明史·刘瑾传》；正德官员奏疏与文集 提供相互检验的起点；但罢黜、逮治和若干诏令可靠；受害人数、每一罪名及其是否全由刘瑾个人决定，不能脱离武宗、内官和文官网络的共同作用。 因此本段仅据材料确认事件的时序、行动和可见后果，并把政策文本与地方实践、报称数字和社会经历分开处理。

\eventanchor{ML-1508-001}{明·1508（正德三年）}
\paragraph{1508（正德三年）：正德朝内廷、人事与巡幸相关文书构成本年中枢政治线索，后出道德化叙事须与原奏疏分开。（材料核查重点：诏令或奏议的形成与发受链）} 这项记录把本章所见的制度运作落实到具体年份。其形成背景是：继承、官僚协调或皇权运作的压力使问题进入朝廷程序。 它带来的、而且仍可由后续材料追踪的改变是：它影响后续人事与政策议程；动机和派系标签不能由单一叙事裁定。 《武宗实录》正德三年条；《明史·本纪》及相关列传；诏令、奏疏或会典版本 提供相互检验的起点；但桥接条目只标示可回查的年度问题与材料链；实录有中央选择性，崇祯期尤其需要奏疏、地方、朝鲜及满文材料交叉。；本条特别核对诏令或奏议的形成与发受链。 因此本段仅据材料确认事件的时序、行动和可见后果，并把政策文本与地方实践、报称数字和社会经历分开处理。

\eventanchor{ML-1508-002}{明·1508（正德三年）}
\paragraph{1508（正德三年）：科举、讲学和官员文集的本年线索可用于追踪知识网络，主张与行政结果须分开。} 这项记录把本章所见的制度运作落实到具体年份。其形成背景是：制度、教育或知识传播的需要推动了相关行动或文本形成。 它带来的、而且仍可由后续材料追踪的改变是：它提供制度和传播史的锚点，不能据此推断社会整体接受。 《武宗实录》正德三年条；《明史·选举志》或《艺文志》；文集、碑刻或书目材料 提供相互检验的起点；但桥接条目只标示可回查的年度问题与材料链；实录有中央选择性，崇祯期尤其需要奏疏、地方、朝鲜及满文材料交叉。 因此本段仅据材料确认事件的时序、行动和可见后果，并把政策文本与地方实践、报称数字和社会经历分开处理。

\eventanchor{ML-1508-003}{明·1508（正德三年）}
\paragraph{1508（正德三年）：嘉靖初继统礼制的本年文书构成大礼议过程锚点，礼制理由和政治效应不可简化为单一动机。} 这项记录把本章所见的制度运作落实到具体年份。其形成背景是：继承、官僚协调或皇权运作的压力使问题进入朝廷程序。 它带来的、而且仍可由后续材料追踪的改变是：它影响后续人事与政策议程；动机和派系标签不能由单一叙事裁定。 《武宗实录》正德三年条；《明史·本纪》及相关列传；诏令、奏疏或会典版本 提供相互检验的起点；但桥接条目只标示可回查的年度问题与材料链；实录有中央选择性，崇祯期尤其需要奏疏、地方、朝鲜及满文材料交叉。 因此本段仅据材料确认事件的时序、行动和可见后果，并把政策文本与地方实践、报称数字和社会经历分开处理。

\eventanchor{ML-1508-004}{明·1508（正德三年）}
\paragraph{1508（正德三年）：正德朝内廷、人事与巡幸相关文书构成本年中枢政治线索，后出道德化叙事须与原奏疏分开。（材料核查重点：报称信息的来源、抄传与行政分类）} 这项记录把本章所见的制度运作落实到具体年份。其形成背景是：继承、官僚协调或皇权运作的压力使问题进入朝廷程序。 它带来的、而且仍可由后续材料追踪的改变是：它影响后续人事与政策议程；动机和派系标签不能由单一叙事裁定。 《武宗实录》正德三年条；《明史·本纪》及相关列传；诏令、奏疏或会典版本 提供相互检验的起点；但桥接条目只标示可回查的年度问题与材料链；实录有中央选择性，崇祯期尤其需要奏疏、地方、朝鲜及满文材料交叉。；本条特别核对报称信息的来源、抄传与行政分类。 因此本段仅据材料确认事件的时序、行动和可见后果，并把政策文本与地方实践、报称数字和社会经历分开处理。

\eventanchor{MM-1508-LIU-JIN-LAND-SURVEY}{明·1508（正德三年）}
\paragraph{1508（正德三年）：刘瑾推动清丈田亩和赋役整饬，以核查隐漏田土} 这项记录把本章所见的制度运作落实到具体年份。其形成背景是：内廷财用和军费需求上升，既有黄册、鱼鳞图册与征收额之间的脱节受到关注。 它带来的、而且仍可由后续材料追踪的改变是：清丈加重地方文书与征解压力，并为刘瑾失势后对其苛政的追责提供了政治材料。 《武宗实录》正德三年户部条；《明史·食货志》；地方赋役册与州县志相关条目 提供相互检验的起点；但清丈命令及其与刘瑾政策相关可靠；各地报亩、增税和实际执行差异极大，不能据诏令写成全国完成。 因此本段仅据材料确认事件的时序、行动和可见后果，并把政策文本与地方实践、报称数字和社会经历分开处理。

\eventanchor{ML-1509-001}{明·1509（正德四年）八月}
\paragraph{1509（正德四年）八月：辽东两驻军叛乱，分白银2500两后平息} 这项记录把本章所见的制度运作落实到具体年份。其形成背景是：朝贡名义、边境安全与港口贸易的利益使交涉持续发生。 它带来的、而且仍可由后续材料追踪的改变是：它为后续册封、互市或海防安排留下文书与跨政权比较线索。 《武宗实录》正德四年条；《明史·外国传》；《朝鲜王朝实录》、琉球《历代宝案》或海外档案 提供相互检验的起点；但译写候选以编年材料定位；实录记载反映朝廷选择，崇祯期须另以奏疏、地方及敌对政权材料交叉。 因此本段仅据材料确认事件的时序、行动和可见后果，并把政策文本与地方实践、报称数字和社会经历分开处理。

\eventanchor{ML-1509-002}{明·1509（正德四年）}
\paragraph{1509（正德四年）：地方反抗、边镇调兵和宗藩危机的本年军政材料标记跨区域动员，军数与战果需要复核。} 这项记录把本章所见的制度运作落实到具体年份。其形成背景是：前期边防、地方武装或跨区域资源调动的压力推动了该行动。 它带来的、而且仍可由后续材料追踪的改变是：它改变了后续调兵、补给或地方秩序的条件；战果和伤亡须分开核对。 《武宗实录》正德四年条；《明史·兵志》及相关列传；边镇、地方志或对方材料 提供相互检验的起点；但桥接条目只标示可回查的年度问题与材料链；实录有中央选择性，崇祯期尤其需要奏疏、地方、朝鲜及满文材料交叉。 因此本段仅据材料确认事件的时序、行动和可见后果，并把政策文本与地方实践、报称数字和社会经历分开处理。

\eventanchor{ML-1509-003}{明·1509（正德四年）}
\paragraph{1509（正德四年）：嘉靖初继统礼制的本年文书构成大礼议过程锚点，礼制理由和政治效应不可简化为单一动机。} 这项记录把本章所见的制度运作落实到具体年份。其形成背景是：继承、官僚协调或皇权运作的压力使问题进入朝廷程序。 它带来的、而且仍可由后续材料追踪的改变是：它影响后续人事与政策议程；动机和派系标签不能由单一叙事裁定。 《武宗实录》正德四年条；《明史·本纪》及相关列传；诏令、奏疏或会典版本 提供相互检验的起点；但桥接条目只标示可回查的年度问题与材料链；实录有中央选择性，崇祯期尤其需要奏疏、地方、朝鲜及满文材料交叉。 因此本段仅据材料确认事件的时序、行动和可见后果，并把政策文本与地方实践、报称数字和社会经历分开处理。

\eventanchor{ML-1509-004}{明·1509（正德四年）}
\paragraph{1509（正德四年）：正德朝内廷、人事与巡幸相关文书构成本年中枢政治线索，后出道德化叙事须与原奏疏分开。} 这项记录把本章所见的制度运作落实到具体年份。其形成背景是：继承、官僚协调或皇权运作的压力使问题进入朝廷程序。 它带来的、而且仍可由后续材料追踪的改变是：它影响后续人事与政策议程；动机和派系标签不能由单一叙事裁定。 《武宗实录》正德四年条；《明史·本纪》及相关列传；诏令、奏疏或会典版本 提供相互检验的起点；但桥接条目只标示可回查的年度问题与材料链；实录有中央选择性，崇祯期尤其需要奏疏、地方、朝鲜及满文材料交叉。 因此本段仅据材料确认事件的时序、行动和可见后果，并把政策文本与地方实践、报称数字和社会经历分开处理。

\eventanchor{ML-1510-001}{明·1510（正德五年）五月12日}
\paragraph{1510（正德五年）五月12日：安化王叛乱：山西朱之璠叛乱} 这项记录把本章所见的制度运作落实到具体年份。其形成背景是：前期边防、地方武装或跨区域资源调动的压力推动了该行动。 它带来的、而且仍可由后续材料追踪的改变是：它改变了后续调兵、补给或地方秩序的条件；战果和伤亡须分开核对。 《武宗实录》正德五年条；《明史·兵志》及相关列传；边镇、地方志或对方材料 提供相互检验的起点；但译写候选以编年材料定位；实录记载反映朝廷选择，崇祯期须另以奏疏、地方及敌对政权材料交叉。 因此本段仅据材料确认事件的时序、行动和可见后果，并把政策文本与地方实践、报称数字和社会经历分开处理。

\eventanchor{ML-1510-002}{明·1510（正德五年）五月30日}
\paragraph{1510（正德五年）五月30日：安化王叛乱：朱之璠被俘} 这项记录把本章所见的制度运作落实到具体年份。其形成背景是：前期边防、地方武装或跨区域资源调动的压力推动了该行动。 它带来的、而且仍可由后续材料追踪的改变是：它改变了后续调兵、补给或地方秩序的条件；战果和伤亡须分开核对。 《武宗实录》正德五年条；《明史·兵志》及相关列传；边镇、地方志或对方材料 提供相互检验的起点；但译写候选以编年材料定位；实录记载反映朝廷选择，崇祯期须另以奏疏、地方及敌对政权材料交叉。 因此本段仅据材料确认事件的时序、行动和可见后果，并把政策文本与地方实践、报称数字和社会经历分开处理。

\eventanchor{ML-1510-003}{明·1510（正德五年）}
\paragraph{1510（正德五年）：达延汗征服鄂尔多斯河套地区} 这项记录把本章所见的制度运作落实到具体年份。其形成背景是：前期边防、地方武装或跨区域资源调动的压力推动了该行动。 它带来的、而且仍可由后续材料追踪的改变是：它改变了后续调兵、补给或地方秩序的条件；战果和伤亡须分开核对。 《武宗实录》正德五年条；《明史·兵志》及相关列传；边镇、地方志或对方材料 提供相互检验的起点；但译写候选以编年材料定位；实录记载反映朝廷选择，崇祯期须另以奏疏、地方及敌对政权材料交叉。 因此本段仅据材料确认事件的时序、行动和可见后果，并把政策文本与地方实践、报称数字和社会经历分开处理。

\eventanchor{MM-1510-ANHUA-REBELLION}{明·1510（正德五年）五月}
\paragraph{1510（正德五年）五月：安化王朱寘鐇在宁夏起兵，旋被当地官军与将领平定} 这项记录把本章所见的制度运作落实到具体年份。其形成背景是：朱寘鐇利用宁夏军镇对刘瑾政策和军饷的不满，试图争取边军支持。 它带来的、而且仍可由后续材料追踪的改变是：叛乱失败削弱了刘瑾的政治位置；宗藩与边镇军事资源的结合再次被朝廷视为风险。 《武宗实录》正德五年五月条；《明史·安化王传》；《明史·仇钺传》 提供相互检验的起点；但起兵和迅速平定可靠；刘瑾是否为叛乱唯一诱因及参与规模不能照录案狱叙事。 因此本段仅据材料确认事件的时序、行动和可见后果，并把政策文本与地方实践、报称数字和社会经历分开处理。

\subsection{1510年前后的续接}

\eventanchor{MM-1510-LIU-JIN-EXECUTION}{明·1510（正德五年）八月}
\paragraph{1510（正德五年）八月：刘瑾被以谋反等罪处死} 这项记录把本章所见的制度运作落实到具体年份。其形成背景是：安化王之变暴露刘瑾的政治脆弱，杨一清、张永等利用边事与宫廷关系促成其失势。 它带来的、而且仍可由后续材料追踪的改变是：刘瑾集团瓦解，正德朝并未由此摆脱宦官和近侍政治；人事、财政政策随之重新调整。 《武宗实录》正德五年八月条；《明史·刘瑾传》；《明史·刑法志》 提供相互检验的起点；但刘瑾被处死及官方罪名可靠；案件罗列的财产、同党和谋反细节是政变后审讯材料，需谨慎使用。 因此本段仅据材料确认事件的时序、行动和可见后果，并把政策文本与地方实践、报称数字和社会经历分开处理。

\eventanchor{ML-1511-001}{明·1511（正德六年）二月}
\paragraph{1511（正德六年）二月：北京周边土匪起义} 这项记录把本章所见的制度运作落实到具体年份。其形成背景是：继承、官僚协调或皇权运作的压力使问题进入朝廷程序。 它带来的、而且仍可由后续材料追踪的改变是：它影响后续人事与政策议程；动机和派系标签不能由单一叙事裁定。 《武宗实录》正德六年条；《明史·本纪》及相关列传；诏令、奏疏或会典版本 提供相互检验的起点；但译写候选以编年材料定位；实录记载反映朝廷选择，崇祯期须另以奏疏、地方及敌对政权材料交叉。 因此本段仅据材料确认事件的时序、行动和可见后果，并把政策文本与地方实践、报称数字和社会经历分开处理。

\eventanchor{ML-1511-002}{明·1511（正德六年）十月}
\paragraph{1511（正德六年）十月：北京各地土匪烧毁皇家粮车} 这项记录把本章所见的制度运作落实到具体年份。其形成背景是：继承、官僚协调或皇权运作的压力使问题进入朝廷程序。 它带来的、而且仍可由后续材料追踪的改变是：它影响后续人事与政策议程；动机和派系标签不能由单一叙事裁定。 《武宗实录》正德六年条；《明史·本纪》及相关列传；诏令、奏疏或会典版本 提供相互检验的起点；但译写候选以编年材料定位；实录记载反映朝廷选择，崇祯期须另以奏疏、地方及敌对政权材料交叉。 因此本段仅据材料确认事件的时序、行动和可见后果，并把政策文本与地方实践、报称数字和社会经历分开处理。

\eventanchor{ML-1511-003}{明·1511（正德六年）}
\paragraph{1511（正德六年）：占领马六甲 (1511)：马六甲苏丹国向葡萄牙人发出求援请求} 这项记录把本章所见的制度运作落实到具体年份。其形成背景是：朝贡名义、边境安全与港口贸易的利益使交涉持续发生。 它带来的、而且仍可由后续材料追踪的改变是：它为后续册封、互市或海防安排留下文书与跨政权比较线索。 《武宗实录》正德六年条；《明史·外国传》；《朝鲜王朝实录》、琉球《历代宝案》或海外档案 提供相互检验的起点；但译写候选以编年材料定位；实录记载反映朝廷选择，崇祯期须另以奏疏、地方及敌对政权材料交叉。 因此本段仅据材料确认事件的时序、行动和可见后果，并把政策文本与地方实践、报称数字和社会经历分开处理。

\eventanchor{MM-1511-LIU-LIU-QI-UPRISING}{明·1511（正德六年）}
\paragraph{1511（正德六年）：刘六、刘七等在山东、北直隶一带聚众活动，成为跨区域治安危机} 这项记录把本章所见的制度运作落实到具体年份。其形成背景是：赋役、灾害与人口流动使部分流民脱离编户控制，运河与平原交通便利了其跨县活动。 它带来的、而且仍可由后续材料追踪的改变是：朝廷调动军队并依靠地方防卫，至1512年前后逐步击散主要集团；沿线社会的损失须以地方材料分别核查。 《武宗实录》正德六年条；《明史·流贼传》；山东、北直隶地方志灾荒与兵事条 提供相互检验的起点；但活动的时间和区域大体可靠；“流贼”称谓遮蔽了参与者的生计、兵役与迁徙背景，人数尤其不可靠。 因此本段仅据材料确认事件的时序、行动和可见后果，并把政策文本与地方实践、报称数字和社会经历分开处理。

\eventanchor{ML-1512-001}{明·1512（正德七年）一月}
\paragraph{1512（正德七年）一月：土匪袭击霸州} 这项记录把本章所见的制度运作落实到具体年份。其形成背景是：继承、官僚协调或皇权运作的压力使问题进入朝廷程序。 它带来的、而且仍可由后续材料追踪的改变是：它影响后续人事与政策议程；动机和派系标签不能由单一叙事裁定。 《武宗实录》正德七年条；《明史·本纪》及相关列传；诏令、奏疏或会典版本 提供相互检验的起点；但译写候选以编年材料定位；实录记载反映朝廷选择，崇祯期须另以奏疏、地方及敌对政权材料交叉。 因此本段仅据材料确认事件的时序、行动和可见后果，并把政策文本与地方实践、报称数字和社会经历分开处理。

\eventanchor{ML-1512-002}{明·1512（正德七年）九月7日}
\paragraph{1512（正德七年）九月7日：强盗大军被击败} 这项记录把本章所见的制度运作落实到具体年份。其形成背景是：前期边防、地方武装或跨区域资源调动的压力推动了该行动。 它带来的、而且仍可由后续材料追踪的改变是：它改变了后续调兵、补给或地方秩序的条件；战果和伤亡须分开核对。 《武宗实录》正德七年条；《明史·兵志》及相关列传；边镇、地方志或对方材料 提供相互检验的起点；但译写候选以编年材料定位；实录记载反映朝廷选择，崇祯期须另以奏疏、地方及敌对政权材料交叉。 因此本段仅据材料确认事件的时序、行动和可见后果，并把政策文本与地方实践、报称数字和社会经历分开处理。

\eventanchor{ML-1512-003}{明·1512（正德七年）}
\paragraph{1512（正德七年）：地方反抗、边镇调兵和宗藩危机的本年军政材料标记跨区域动员，军数与战果需要复核。} 这项记录把本章所见的制度运作落实到具体年份。其形成背景是：前期边防、地方武装或跨区域资源调动的压力推动了该行动。 它带来的、而且仍可由后续材料追踪的改变是：它改变了后续调兵、补给或地方秩序的条件；战果和伤亡须分开核对。 《武宗实录》正德七年条；《明史·兵志》及相关列传；边镇、地方志或对方材料 提供相互检验的起点；但桥接条目只标示可回查的年度问题与材料链；实录有中央选择性，崇祯期尤其需要奏疏、地方、朝鲜及满文材料交叉。 因此本段仅据材料确认事件的时序、行动和可见后果，并把政策文本与地方实践、报称数字和社会经历分开处理。

\eventanchor{MM-1512-LIU-LIU-QI-SUPPRESSION}{明·1512（正德七年）}
\paragraph{1512（正德七年）：官军在山东、北直隶分路追击刘六、刘七等集团，主要武装至年内瓦解} 这项记录把本章所见的制度运作落实到具体年份。其形成背景是：1511年起跨县流动的武装借助运河与平原交通维持补给，朝廷调兵和地方防守逐渐压缩其活动空间。 它带来的、而且仍可由后续材料追踪的改变是：主要首领被杀或失散后危机降温，但流民、灾荒和基层赋役问题未被追剿解决，沿线治安仍需地方社会承担。 《武宗实录》正德七年山东与北直隶条；《明史·流贼传》；德州、临清等地地方志兵事条 提供相互检验的起点；但追剿、首领结局和朝廷奖惩大体可靠；官方所称歼获、余众去向和沿线受害程度主要来自战报，须以州县材料复核。 因此本段仅据材料确认事件的时序、行动和可见后果，并把政策文本与地方实践、报称数字和社会经历分开处理。

\eventanchor{ML-1513-001}{明·1513（正德八年）}
\paragraph{1513（正德八年）：清丈、库藏、差役与征解的本年记录显示财政监管压力，整饬命令不等于隐漏已被清除。（材料核查重点：诏令或奏议的形成与发受链）} 这项记录把本章所见的制度运作落实到具体年份。其形成背景是：财政征解、交通工程或货币支付的压力使相关措施进入政务。 它带来的、而且仍可由后续材料追踪的改变是：其法定安排不等于地方执行，后续须以地域材料追踪负担与成效。 《武宗实录》正德八年条；《明会典》相关条目；地方志、黄册/契约/河工材料 提供相互检验的起点；但桥接条目只标示可回查的年度问题与材料链；实录有中央选择性，崇祯期尤其需要奏疏、地方、朝鲜及满文材料交叉。；本条特别核对诏令或奏议的形成与发受链。 因此本段仅据材料确认事件的时序、行动和可见后果，并把政策文本与地方实践、报称数字和社会经历分开处理。

\eventanchor{ML-1513-002}{明·1513（正德八年）}
\paragraph{1513（正德八年）：正德朝内廷、人事与巡幸相关文书构成本年中枢政治线索，后出道德化叙事须与原奏疏分开。} 这项记录把本章所见的制度运作落实到具体年份。其形成背景是：继承、官僚协调或皇权运作的压力使问题进入朝廷程序。 它带来的、而且仍可由后续材料追踪的改变是：它影响后续人事与政策议程；动机和派系标签不能由单一叙事裁定。 《武宗实录》正德八年条；《明史·本纪》及相关列传；诏令、奏疏或会典版本 提供相互检验的起点；但桥接条目只标示可回查的年度问题与材料链；实录有中央选择性，崇祯期尤其需要奏疏、地方、朝鲜及满文材料交叉。 因此本段仅据材料确认事件的时序、行动和可见后果，并把政策文本与地方实践、报称数字和社会经历分开处理。

\eventanchor{ML-1513-003}{明·1513（正德八年）}
\paragraph{1513（正德八年）：清丈、库藏、差役与征解的本年记录显示财政监管压力，整饬命令不等于隐漏已被清除。（材料核查重点：报称信息的来源、抄传与行政分类）} 这项记录把本章所见的制度运作落实到具体年份。其形成背景是：财政征解、交通工程或货币支付的压力使相关措施进入政务。 它带来的、而且仍可由后续材料追踪的改变是：其法定安排不等于地方执行，后续须以地域材料追踪负担与成效。 《武宗实录》正德八年条；《明会典》相关条目；地方志、黄册/契约/河工材料 提供相互检验的起点；但桥接条目只标示可回查的年度问题与材料链；实录有中央选择性，崇祯期尤其需要奏疏、地方、朝鲜及满文材料交叉。；本条特别核对报称信息的来源、抄传与行政分类。 因此本段仅据材料确认事件的时序、行动和可见后果，并把政策文本与地方实践、报称数字和社会经历分开处理。

\eventanchor{MM-1513-PORTUGUESE-TAMAO-CONTACT}{明·1513（正德八年）}
\paragraph{1513（正德八年）：葡萄牙航海者在屯门一带停泊并尝试贸易，为后续使团接触提供海上先例} 这项记录把本章所见的制度运作落实到具体年份。其形成背景是：葡萄牙人进入印度洋—马六甲航路后寻求华南货源，而珠江口已有商人、翻译与官府共同维系的港口准入规则。 它带来的、而且仍可由后续材料追踪的改变是：此类非正式接触未取得稳定外交资格，却使中介与地方官更早面对欧洲船只；1517年使团入粤应置于此背景。 若昂·德·巴罗斯《亚洲》早期航行记述；《明史·佛郎机传》；珠江口港航研究 提供相互检验的起点；但早期抵达主要见于葡萄牙航海叙事，汉文记录对人名和地点较少；屯门是否等同葡文所称Tamão、到达月日与交易规模均有争论。 因此本段仅据材料确认事件的时序、行动和可见后果，并把政策文本与地方实践、报称数字和社会经历分开处理。

\eventanchor{ML-1514-001}{明·1514（正德九年）二月10日}
\paragraph{1514（正德九年）二月10日：宫殿庭院内的火药帐篷着火，烧毁了住宅宫殿} 这项记录把本章所见的制度运作落实到具体年份。其形成背景是：自然风险与地方报告促成朝廷或地方的应对记录。 它带来的、而且仍可由后续材料追踪的改变是：赈济、迁徙和损失的实际范围仍须以受灾地区材料分别核验。 《武宗实录》正德九年条；《明史·五行志》；受灾府县地方志与奏疏 提供相互检验的起点；但译写候选以编年材料定位；实录记载反映朝廷选择，崇祯期须另以奏疏、地方及敌对政权材料交叉。 因此本段仅据材料确认事件的时序、行动和可见后果，并把政策文本与地方实践、报称数字和社会经历分开处理。

\eventanchor{ML-1514-002}{明·1514（正德九年）九月}
\paragraph{1514（正德九年）九月：正德皇帝被老虎咬成重伤} 这项记录把本章所见的制度运作落实到具体年份。其形成背景是：继承、官僚协调或皇权运作的压力使问题进入朝廷程序。 它带来的、而且仍可由后续材料追踪的改变是：它影响后续人事与政策议程；动机和派系标签不能由单一叙事裁定。 《武宗实录》正德九年条；《明史·本纪》及相关列传；诏令、奏疏或会典版本 提供相互检验的起点；但译写候选以编年材料定位；实录记载反映朝廷选择，崇祯期须另以奏疏、地方及敌对政权材料交叉。 因此本段仅据材料确认事件的时序、行动和可见后果，并把政策文本与地方实践、报称数字和社会经历分开处理。

\eventanchor{ML-1514-003}{明·1514（正德九年）}
\paragraph{1514（正德九年）：东南港口、日本使团或葡萄牙接触的本年材料为海上交往留档，朝贡、私贸与冲突不可混称。} 这项记录把本章所见的制度运作落实到具体年份。其形成背景是：朝贡名义、边境安全与港口贸易的利益使交涉持续发生。 它带来的、而且仍可由后续材料追踪的改变是：它为后续册封、互市或海防安排留下文书与跨政权比较线索。 《武宗实录》正德九年条；《明史·外国传》；《朝鲜王朝实录》、琉球《历代宝案》或海外档案 提供相互检验的起点；但桥接条目只标示可回查的年度问题与材料链；实录有中央选择性，崇祯期尤其需要奏疏、地方、朝鲜及满文材料交叉。 因此本段仅据材料确认事件的时序、行动和可见后果，并把政策文本与地方实践、报称数字和社会经历分开处理。

\eventanchor{MM-1514-TURPAN-HAMI-ANNEXATION}{明·1514（正德九年）}
\paragraph{1514（正德九年）：吐鲁番势力再次夺取哈密，明朝维系该绿洲朝贡政权的能力显著下降} 这项记录把本章所见的制度运作落实到具体年份。其形成背景是：弘治朝后的哈密继承安排缺乏稳定军力保障，吐鲁番可凭更近的补给和政治联盟重新施压。 它带来的、而且仍可由后续材料追踪的改变是：明廷对哈密的直接扶持更难持续，嘉峪关外的册封体系和关市管理转而面对吐鲁番主导的现实。 《武宗实录》正德九年哈密条；《明史·西域传》；吐鲁番文书与中亚编年材料 提供相互检验的起点；但夺取及明廷获报可靠；是否构成完全且永久的行政并入、当地社群反应和军队数额在汉文与中亚材料中不能简单等同。 因此本段仅据材料确认事件的时序、行动和可见后果，并把政策文本与地方实践、报称数字和社会经历分开处理。

\eventanchor{ML-1515-001}{明·1515（正德十年）夏}
\paragraph{1515（正德十年）夏：京师、侍卫三万出动，重建宫殿} 这项记录把本章所见的制度运作落实到具体年份。其形成背景是：继承、官僚协调或皇权运作的压力使问题进入朝廷程序。 它带来的、而且仍可由后续材料追踪的改变是：它影响后续人事与政策议程；动机和派系标签不能由单一叙事裁定。 《武宗实录》正德十年条；《明史·本纪》及相关列传；诏令、奏疏或会典版本 提供相互检验的起点；但译写候选以编年材料定位；实录记载反映朝廷选择，崇祯期须另以奏疏、地方及敌对政权材料交叉。 因此本段仅据材料确认事件的时序、行动和可见后果，并把政策文本与地方实践、报称数字和社会经历分开处理。

\eventanchor{ML-1515-002}{明·1515（正德十年）}
\paragraph{1515（正德十年）：东南港口、日本使团或葡萄牙接触的本年材料为海上交往留档，朝贡、私贸与冲突不可混称。} 这项记录把本章所见的制度运作落实到具体年份。其形成背景是：朝贡名义、边境安全与港口贸易的利益使交涉持续发生。 它带来的、而且仍可由后续材料追踪的改变是：它为后续册封、互市或海防安排留下文书与跨政权比较线索。 《武宗实录》正德十年条；《明史·外国传》；《朝鲜王朝实录》、琉球《历代宝案》或海外档案 提供相互检验的起点；但桥接条目只标示可回查的年度问题与材料链；实录有中央选择性，崇祯期尤其需要奏疏、地方、朝鲜及满文材料交叉。 因此本段仅据材料确认事件的时序、行动和可见后果，并把政策文本与地方实践、报称数字和社会经历分开处理。

\eventanchor{ML-1515-003}{明·1515（正德十年）}
\paragraph{1515（正德十年）：地方反抗、边镇调兵和宗藩危机的本年军政材料标记跨区域动员，军数与战果需要复核。} 这项记录把本章所见的制度运作落实到具体年份。其形成背景是：前期边防、地方武装或跨区域资源调动的压力推动了该行动。 它带来的、而且仍可由后续材料追踪的改变是：它改变了后续调兵、补给或地方秩序的条件；战果和伤亡须分开核对。 《武宗实录》正德十年条；《明史·兵志》及相关列传；边镇、地方志或对方材料 提供相互检验的起点；但桥接条目只标示可回查的年度问题与材料链；实录有中央选择性，崇祯期尤其需要奏疏、地方、朝鲜及满文材料交叉。 因此本段仅据材料确认事件的时序、行动和可见后果，并把政策文本与地方实践、报称数字和社会经历分开处理。

\subsection{1515年前后的续接}

\eventanchor{ML-1515-004}{明·1515（正德十年）}
\paragraph{1515（正德十年）：科举、讲学和官员文集的本年线索可用于追踪知识网络，主张与行政结果须分开。} 这项记录把本章所见的制度运作落实到具体年份。其形成背景是：制度、教育或知识传播的需要推动了相关行动或文本形成。 它带来的、而且仍可由后续材料追踪的改变是：它提供制度和传播史的锚点，不能据此推断社会整体接受。 《武宗实录》正德十年条；《明史·选举志》或《艺文志》；文集、碑刻或书目材料 提供相互检验的起点；但桥接条目只标示可回查的年度问题与材料链；实录有中央选择性，崇祯期尤其需要奏疏、地方、朝鲜及满文材料交叉。 因此本段仅据材料确认事件的时序、行动和可见后果，并把政策文本与地方实践、报称数字和社会经历分开处理。

\eventanchor{ML-1516-001}{明·1516（正德十一年）}
\paragraph{1516（正德十一年）：地方反抗、边镇调兵和宗藩危机的本年军政材料标记跨区域动员，军数与战果需要复核。} 这项记录把本章所见的制度运作落实到具体年份。其形成背景是：前期边防、地方武装或跨区域资源调动的压力推动了该行动。 它带来的、而且仍可由后续材料追踪的改变是：它改变了后续调兵、补给或地方秩序的条件；战果和伤亡须分开核对。 《武宗实录》正德十一年条；《明史·兵志》及相关列传；边镇、地方志或对方材料 提供相互检验的起点；但桥接条目只标示可回查的年度问题与材料链；实录有中央选择性，崇祯期尤其需要奏疏、地方、朝鲜及满文材料交叉。 因此本段仅据材料确认事件的时序、行动和可见后果，并把政策文本与地方实践、报称数字和社会经历分开处理。

\eventanchor{ML-1516-002}{明·1516（正德十一年）}
\paragraph{1516（正德十一年）：科举、讲学和官员文集的本年线索可用于追踪知识网络，主张与行政结果须分开。} 这项记录把本章所见的制度运作落实到具体年份。其形成背景是：制度、教育或知识传播的需要推动了相关行动或文本形成。 它带来的、而且仍可由后续材料追踪的改变是：它提供制度和传播史的锚点，不能据此推断社会整体接受。 《武宗实录》正德十一年条；《明史·选举志》或《艺文志》；文集、碑刻或书目材料 提供相互检验的起点；但桥接条目只标示可回查的年度问题与材料链；实录有中央选择性，崇祯期尤其需要奏疏、地方、朝鲜及满文材料交叉。 因此本段仅据材料确认事件的时序、行动和可见后果，并把政策文本与地方实践、报称数字和社会经历分开处理。

\eventanchor{ML-1516-003}{明·1516（正德十一年）}
\paragraph{1516（正德十一年）：清丈、库藏、差役与征解的本年记录显示财政监管压力，整饬命令不等于隐漏已被清除。} 这项记录把本章所见的制度运作落实到具体年份。其形成背景是：财政征解、交通工程或货币支付的压力使相关措施进入政务。 它带来的、而且仍可由后续材料追踪的改变是：其法定安排不等于地方执行，后续须以地域材料追踪负担与成效。 《武宗实录》正德十一年条；《明会典》相关条目；地方志、黄册/契约/河工材料 提供相互检验的起点；但桥接条目只标示可回查的年度问题与材料链；实录有中央选择性，崇祯期尤其需要奏疏、地方、朝鲜及满文材料交叉。 因此本段仅据材料确认事件的时序、行动和可见后果，并把政策文本与地方实践、报称数字和社会经历分开处理。

\eventanchor{ML-1516-004}{明·1516（正德十一年）}
\paragraph{1516（正德十一年）：嘉靖初继统礼制的本年文书构成大礼议过程锚点，礼制理由和政治效应不可简化为单一动机。} 这项记录把本章所见的制度运作落实到具体年份。其形成背景是：继承、官僚协调或皇权运作的压力使问题进入朝廷程序。 它带来的、而且仍可由后续材料追踪的改变是：它影响后续人事与政策议程；动机和派系标签不能由单一叙事裁定。 《武宗实录》正德十一年条；《明史·本纪》及相关列传；诏令、奏疏或会典版本 提供相互检验的起点；但桥接条目只标示可回查的年度问题与材料链；实录有中央选择性，崇祯期尤其需要奏疏、地方、朝鲜及满文材料交叉。 因此本段仅据材料确认事件的时序、行动和可见后果，并把政策文本与地方实践、报称数字和社会经历分开处理。

\eventanchor{ML-1517-001}{明·1517（正德十二年）十月16日}
\paragraph{1517（正德十二年）十月16日：达延汗袭击明朝} 这项记录把本章所见的制度运作落实到具体年份。其形成背景是：前期边防、地方武装或跨区域资源调动的压力推动了该行动。 它带来的、而且仍可由后续材料追踪的改变是：它改变了后续调兵、补给或地方秩序的条件；战果和伤亡须分开核对。 《武宗实录》正德十二年条；《明史·兵志》及相关列传；边镇、地方志或对方材料 提供相互检验的起点；但译写候选以编年材料定位；实录记载反映朝廷选择，崇祯期须另以奏疏、地方及敌对政权材料交叉。 因此本段仅据材料确认事件的时序、行动和可见后果，并把政策文本与地方实践、报称数字和社会经历分开处理。

\eventanchor{ML-1517-002}{明·1517（正德十二年）十月20日}
\paragraph{1517（正德十二年）十月20日：正德皇帝击退达延汗的袭击} 这项记录把本章所见的制度运作落实到具体年份。其形成背景是：前期边防、地方武装或跨区域资源调动的压力推动了该行动。 它带来的、而且仍可由后续材料追踪的改变是：它改变了后续调兵、补给或地方秩序的条件；战果和伤亡须分开核对。 《武宗实录》正德十二年条；《明史·兵志》及相关列传；边镇、地方志或对方材料 提供相互检验的起点；但译写候选以编年材料定位；实录记载反映朝廷选择，崇祯期须另以奏疏、地方及敌对政权材料交叉。 因此本段仅据材料确认事件的时序、行动和可见后果，并把政策文本与地方实践、报称数字和社会经历分开处理。

\eventanchor{ML-1517-003}{明·1517（正德十二年）}
\paragraph{1517（正德十二年）：托梅·皮雷抵达广州} 这项记录把本章所见的制度运作落实到具体年份。其形成背景是：继承、官僚协调或皇权运作的压力使问题进入朝廷程序。 它带来的、而且仍可由后续材料追踪的改变是：它影响后续人事与政策议程；动机和派系标签不能由单一叙事裁定。 《武宗实录》正德十二年条；《明史·本纪》及相关列传；诏令、奏疏或会典版本 提供相互检验的起点；但译写候选以编年材料定位；实录记载反映朝廷选择，崇祯期须另以奏疏、地方及敌对政权材料交叉。 因此本段仅据材料确认事件的时序、行动和可见后果，并把政策文本与地方实践、报称数字和社会经历分开处理。

\eventanchor{MM-1517-PORTUGUESE-EMBASSY}{明·1517（正德十二年）}
\paragraph{1517（正德十二年）：葡萄牙使团抵达广东，托名满剌加使者请求入朝} 这项记录把本章所见的制度运作落实到具体年份。其形成背景是：葡萄牙人已进入亚洲海上贸易网络，试图以熟悉的朝贡语言取得在华贸易和外交准入。 它带来的、而且仍可由后续材料追踪的改变是：广东官府开始面对新的海上来者；满剌加问题和港口秩序使使团最终卷入更深的外交冲突。 《武宗实录》正德十二年广东条；《明史·佛郎机传》；葡萄牙使团书信与航海记载 提供相互检验的起点；但使团抵粤和请求入朝可由中葡材料互证；使团身份、翻译过程及其代表的贸易利益在两类材料中表述不同。 因此本段仅据材料确认事件的时序、行动和可见后果，并把政策文本与地方实践、报称数字和社会经历分开处理。

\eventanchor{ML-1518-001}{明·1518（正德十三年）一月}
\paragraph{1518（正德十三年）一月：正德皇帝因没有给他足够的钱财而被囚禁在北京} 这项记录把本章所见的制度运作落实到具体年份。其形成背景是：继承、官僚协调或皇权运作的压力使问题进入朝廷程序。 它带来的、而且仍可由后续材料追踪的改变是：它影响后续人事与政策议程；动机和派系标签不能由单一叙事裁定。 《武宗实录》正德十三年条；《明史·本纪》及相关列传；诏令、奏疏或会典版本 提供相互检验的起点；但译写候选以编年材料定位；实录记载反映朝廷选择，崇祯期须另以奏疏、地方及敌对政权材料交叉。 因此本段仅据材料确认事件的时序、行动和可见后果，并把政策文本与地方实践、报称数字和社会经历分开处理。

\eventanchor{ML-1518-002}{明·1518（正德十三年）}
\paragraph{1518（正德十三年）：东南港口、日本使团或葡萄牙接触的本年材料为海上交往留档，朝贡、私贸与冲突不可混称。} 这项记录把本章所见的制度运作落实到具体年份。其形成背景是：朝贡名义、边境安全与港口贸易的利益使交涉持续发生。 它带来的、而且仍可由后续材料追踪的改变是：它为后续册封、互市或海防安排留下文书与跨政权比较线索。 《武宗实录》正德十三年条；《明史·外国传》；《朝鲜王朝实录》、琉球《历代宝案》或海外档案 提供相互检验的起点；但桥接条目只标示可回查的年度问题与材料链；实录有中央选择性，崇祯期尤其需要奏疏、地方、朝鲜及满文材料交叉。 因此本段仅据材料确认事件的时序、行动和可见后果，并把政策文本与地方实践、报称数字和社会经历分开处理。

\eventanchor{ML-1518-003}{明·1518（正德十三年）}
\paragraph{1518（正德十三年）：正德朝内廷、人事与巡幸相关文书构成本年中枢政治线索，后出道德化叙事须与原奏疏分开。} 这项记录把本章所见的制度运作落实到具体年份。其形成背景是：继承、官僚协调或皇权运作的压力使问题进入朝廷程序。 它带来的、而且仍可由后续材料追踪的改变是：它影响后续人事与政策议程；动机和派系标签不能由单一叙事裁定。 《武宗实录》正德十三年条；《明史·本纪》及相关列传；诏令、奏疏或会典版本 提供相互检验的起点；但桥接条目只标示可回查的年度问题与材料链；实录有中央选择性，崇祯期尤其需要奏疏、地方、朝鲜及满文材料交叉。 因此本段仅据材料确认事件的时序、行动和可见后果，并把政策文本与地方实践、报称数字和社会经历分开处理。

\eventanchor{MM-1518-WANG-YANGMING-GANZHOU}{明·1518（正德十三年）}
\paragraph{1518（正德十三年）：王守仁在赣南主持军政，推行保甲、招抚与军事征剿并用} 这项记录把本章所见的制度运作落实到具体年份。其形成背景是：赣南、汀漳山区长期存在流民、矿徒与地方武装，常规州县与卫所难以有效控制。 它带来的、而且仍可由后续材料追踪的改变是：招抚、保甲和军事行动结合，形成可在宁王之变中调用的地方军政网络；其强制成本和地方接受度仍有待个案研究。 《武宗实录》正德十三年赣南条；《明史·王守仁传》；《王阳明全集》奏疏与地方志 提供相互检验的起点；但王守仁的任命和主要措施可靠；“良知教化”对基层的实际作用不能由其文集直接推定。 因此本段仅据材料确认事件的时序、行动和可见后果，并把政策文本与地方实践、报称数字和社会经历分开处理。

\eventanchor{ML-1519-001}{明·1519（正德十四年）七月9日}
\paragraph{1519（正德十四年）七月9日：宁王叛乱：江西朱辰灏叛乱} 这项记录把本章所见的制度运作落实到具体年份。其形成背景是：前期边防、地方武装或跨区域资源调动的压力推动了该行动。 它带来的、而且仍可由后续材料追踪的改变是：它改变了后续调兵、补给或地方秩序的条件；战果和伤亡须分开核对。 《武宗实录》正德十四年条；《明史·兵志》及相关列传；边镇、地方志或对方材料 提供相互检验的起点；但译写候选以编年材料定位；实录记载反映朝廷选择，崇祯期须另以奏疏、地方及敌对政权材料交叉。 因此本段仅据材料确认事件的时序、行动和可见后果，并把政策文本与地方实践、报称数字和社会经历分开处理。

\eventanchor{ML-1519-002}{明·1519（正德十四年）七月13日}
\paragraph{1519（正德十四年）七月13日：宁王叛乱：叛军攻克九江} 这项记录把本章所见的制度运作落实到具体年份。其形成背景是：前期边防、地方武装或跨区域资源调动的压力推动了该行动。 它带来的、而且仍可由后续材料追踪的改变是：它改变了后续调兵、补给或地方秩序的条件；战果和伤亡须分开核对。 《武宗实录》正德十四年条；《明史·兵志》及相关列传；边镇、地方志或对方材料 提供相互检验的起点；但译写候选以编年材料定位；实录记载反映朝廷选择，崇祯期须另以奏疏、地方及敌对政权材料交叉。 因此本段仅据材料确认事件的时序、行动和可见后果，并把政策文本与地方实践、报称数字和社会经历分开处理。

\eventanchor{ML-1519-003}{明·1519（正德十四年）七月23日}
\paragraph{1519（正德十四年）七月23日：宁王叛乱：叛军围攻安庆} 这项记录把本章所见的制度运作落实到具体年份。其形成背景是：前期边防、地方武装或跨区域资源调动的压力推动了该行动。 它带来的、而且仍可由后续材料追踪的改变是：它改变了后续调兵、补给或地方秩序的条件；战果和伤亡须分开核对。 《武宗实录》正德十四年条；《明史·兵志》及相关列传；边镇、地方志或对方材料 提供相互检验的起点；但译写候选以编年材料定位；实录记载反映朝廷选择，崇祯期须另以奏疏、地方及敌对政权材料交叉。 因此本段仅据材料确认事件的时序、行动和可见后果，并把政策文本与地方实践、报称数字和社会经历分开处理。

\eventanchor{ML-1519-004}{明·1519（正德十四年）八月9日}
\paragraph{1519（正德十四年）八月9日：宁王叛乱：叛军解围安庆} 这项记录把本章所见的制度运作落实到具体年份。其形成背景是：前期边防、地方武装或跨区域资源调动的压力推动了该行动。 它带来的、而且仍可由后续材料追踪的改变是：它改变了后续调兵、补给或地方秩序的条件；战果和伤亡须分开核对。 《武宗实录》正德十四年条；《明史·兵志》及相关列传；边镇、地方志或对方材料 提供相互检验的起点；但译写候选以编年材料定位；实录记载反映朝廷选择，崇祯期须另以奏疏、地方及敌对政权材料交叉。 因此本段仅据材料确认事件的时序、行动和可见后果，并把政策文本与地方实践、报称数字和社会经历分开处理。

\eventanchor{ML-1519-005}{明·1519（正德十四年）八月13日}
\paragraph{1519（正德十四年）八月13日：宁王叛乱：清军攻克南昌} 这项记录把本章所见的制度运作落实到具体年份。其形成背景是：前期边防、地方武装或跨区域资源调动的压力推动了该行动。 它带来的、而且仍可由后续材料追踪的改变是：它改变了后续调兵、补给或地方秩序的条件；战果和伤亡须分开核对。 《武宗实录》正德十四年条；《明史·兵志》及相关列传；边镇、地方志或对方材料 提供相互检验的起点；但译写候选以编年材料定位；实录记载反映朝廷选择，崇祯期须另以奏疏、地方及敌对政权材料交叉。 因此本段仅据材料确认事件的时序、行动和可见后果，并把政策文本与地方实践、报称数字和社会经历分开处理。

\eventanchor{ML-1519-006}{明·1519（正德十四年）八月15日}
\paragraph{1519（正德十四年）八月15日：宁王叛乱：朱辰灏军大败} 这项记录把本章所见的制度运作落实到具体年份。其形成背景是：前期边防、地方武装或跨区域资源调动的压力推动了该行动。 它带来的、而且仍可由后续材料追踪的改变是：它改变了后续调兵、补给或地方秩序的条件；战果和伤亡须分开核对。 《武宗实录》正德十四年条；《明史·兵志》及相关列传；边镇、地方志或对方材料 提供相互检验的起点；但译写候选以编年材料定位；实录记载反映朝廷选择，崇祯期须另以奏疏、地方及敌对政权材料交叉。 因此本段仅据材料确认事件的时序、行动和可见后果，并把政策文本与地方实践、报称数字和社会经历分开处理。

\eventanchor{ML-1519-007}{明·1519（正德十四年）八月20日}
\paragraph{1519（正德十四年）八月20日：宁王叛乱：朱辰浩逃离舰队被俘} 这项记录把本章所见的制度运作落实到具体年份。其形成背景是：前期边防、地方武装或跨区域资源调动的压力推动了该行动。 它带来的、而且仍可由后续材料追踪的改变是：它改变了后续调兵、补给或地方秩序的条件；战果和伤亡须分开核对。 《武宗实录》正德十四年条；《明史·兵志》及相关列传；边镇、地方志或对方材料 提供相互检验的起点；但译写候选以编年材料定位；实录记载反映朝廷选择，崇祯期须另以奏疏、地方及敌对政权材料交叉。 因此本段仅据材料确认事件的时序、行动和可见后果，并把政策文本与地方实践、报称数字和社会经历分开处理。

\subsection{1519年前后的续接}

\eventanchor{MM-1519-NING-REBELLION}{明·1519（正德十四年）六月}
\paragraph{1519（正德十四年）六月：宁王朱宸濠在南昌起兵，试图沿江扩张} 这项记录把本章所见的制度运作落实到具体年份。其形成背景是：宗藩限制虽已趋严，朱宸濠仍在江西经营人脉和军事资源；中央又受武宗南巡与内廷政治牵制。 它带来的、而且仍可由后续材料追踪的改变是：叛乱将江西军政快速推入战争状态，王守仁的地方部署成为阻止其沿江扩张的关键。 《武宗实录》正德十四年六月条；《明史·宁王宸濠传》；《王阳明全集》平叛奏疏 提供相互检验的起点；但起兵日期、据点和主要行动可靠；其称帝计划、同谋网络和动机大量出自平叛后的案卷，不能无保留采信。 因此本段仅据材料确认事件的时序、行动和可见后果，并把政策文本与地方实践、报称数字和社会经历分开处理。

\eventanchor{MM-1519-WANG-YANGMING-VICTORY}{明·1519（正德十四年）七月}
\paragraph{1519（正德十四年）七月：王守仁率地方军在鄱阳湖、南昌一带击败并俘获朱宸濠} 这项记录把本章所见的制度运作落实到具体年份。其形成背景是：王守仁先前在赣南形成的动员与指挥网络，使其能在中央援军到达前集中兵力。 它带来的、而且仍可由后续材料追踪的改变是：朱宸濠被俘，叛乱主体瓦解；平叛功劳和武宗南巡中的献俘安排则引出新的宫廷摩擦。 《武宗实录》正德十四年七月条；《明史·王守仁传》；《王阳明全集》捷报与奏疏 提供相互检验的起点；但战败和俘获可靠；具体战斗地点、兵数和王守仁独自决定胜负的说法须与战报修辞相区分。 因此本段仅据材料确认事件的时序、行动和可见后果，并把政策文本与地方实践、报称数字和社会经历分开处理。

\eventanchor{MM-1519-WUZONG-SOUTHERN-TOUR}{明·1519（正德十四年）}
\paragraph{1519（正德十四年）：宁王之变已受控制后，武宗仍南巡并以亲临军务为名驻跸江南} 这项记录把本章所见的制度运作落实到具体年份。其形成背景是：武宗希望直接掌握已被平定的宁王案件和南方军事象征，内廷近侍也可借巡幸影响人事、供给和信息。 它带来的、而且仍可由后续材料追踪的改变是：江南州县承担行在供应与警卫压力，王守仁的平叛功劳及献俘安排被重新置入宫廷竞争；巡幸延续至1521年。 《武宗实录》正德十四年南巡条；《明史·武宗本纪》；王守仁与江南官员奏疏 提供相互检验的起点；但南巡、驻跸与随行机构可由实录和官员文集互证；皇帝动机、行程细节与地方耗费在批评奏疏和官修记录中立场差异很大。 因此本段仅据材料确认事件的时序、行动和可见后果，并把政策文本与地方实践、报称数字和社会经历分开处理。

\eventanchor{ML-1520-001}{明·1520（正德十五年）一月}
\paragraph{1520（正德十五年）一月：正德皇帝禁止宰猪} 这项记录把本章所见的制度运作落实到具体年份。其形成背景是：继承、官僚协调或皇权运作的压力使问题进入朝廷程序。 它带来的、而且仍可由后续材料追踪的改变是：它影响后续人事与政策议程；动机和派系标签不能由单一叙事裁定。 《武宗实录》正德十五年条；《明史·本纪》及相关列传；诏令、奏疏或会典版本 提供相互检验的起点；但译写候选以编年材料定位；实录记载反映朝廷选择，崇祯期须另以奏疏、地方及敌对政权材料交叉。 因此本段仅据材料确认事件的时序、行动和可见后果，并把政策文本与地方实践、报称数字和社会经历分开处理。

\eventanchor{ML-1520-002}{明·1520（正德十五年）五月}
\paragraph{1520（正德十五年）五月：葡萄牙人贿赂了广州的一名太监，让他们通过，托梅·皮雷一行抵达南京} 这项记录把本章所见的制度运作落实到具体年份。其形成背景是：朝贡名义、边境安全与港口贸易的利益使交涉持续发生。 它带来的、而且仍可由后续材料追踪的改变是：它为后续册封、互市或海防安排留下文书与跨政权比较线索。 《武宗实录》正德十五年条；《明史·外国传》；《朝鲜王朝实录》、琉球《历代宝案》或海外档案 提供相互检验的起点；但译写候选以编年材料定位；实录记载反映朝廷选择，崇祯期须另以奏疏、地方及敌对政权材料交叉。 因此本段仅据材料确认事件的时序、行动和可见后果，并把政策文本与地方实践、报称数字和社会经历分开处理。

\eventanchor{ML-1520-003}{明·1520（正德十五年）}
\paragraph{1520（正德十五年）：科举、讲学和官员文集的本年线索可用于追踪知识网络，主张与行政结果须分开。} 这项记录把本章所见的制度运作落实到具体年份。其形成背景是：制度、教育或知识传播的需要推动了相关行动或文本形成。 它带来的、而且仍可由后续材料追踪的改变是：它提供制度和传播史的锚点，不能据此推断社会整体接受。 《武宗实录》正德十五年条；《明史·选举志》或《艺文志》；文集、碑刻或书目材料 提供相互检验的起点；但桥接条目只标示可回查的年度问题与材料链；实录有中央选择性，崇祯期尤其需要奏疏、地方、朝鲜及满文材料交叉。 因此本段仅据材料确认事件的时序、行动和可见后果，并把政策文本与地方实践、报称数字和社会经历分开处理。

\eventanchor{ML-1520-004}{明·1520（正德十五年）}
\paragraph{1520（正德十五年）：地方反抗、边镇调兵和宗藩危机的本年军政材料标记跨区域动员，军数与战果需要复核。} 这项记录把本章所见的制度运作落实到具体年份。其形成背景是：前期边防、地方武装或跨区域资源调动的压力推动了该行动。 它带来的、而且仍可由后续材料追踪的改变是：它改变了后续调兵、补给或地方秩序的条件；战果和伤亡须分开核对。 《武宗实录》正德十五年条；《明史·兵志》及相关列传；边镇、地方志或对方材料 提供相互检验的起点；但桥接条目只标示可回查的年度问题与材料链；实录有中央选择性，崇祯期尤其需要奏疏、地方、朝鲜及满文材料交叉。 因此本段仅据材料确认事件的时序、行动和可见后果，并把政策文本与地方实践、报称数字和社会经历分开处理。

\eventanchor{ML-1521-001}{明·1521（正德十六年）四月20日}
\paragraph{1521（正德十六年）四月20日：正德皇帝驾崩} 这项记录把本章所见的制度运作落实到具体年份。其形成背景是：继承、官僚协调或皇权运作的压力使问题进入朝廷程序。 它带来的、而且仍可由后续材料追踪的改变是：它影响后续人事与政策议程；动机和派系标签不能由单一叙事裁定。 《武宗实录》正德十六年条；《明史·本纪》及相关列传；诏令、奏疏或会典版本 提供相互检验的起点；但译写候选以编年材料定位；实录记载反映朝廷选择，崇祯期须另以奏疏、地方及敌对政权材料交叉。 因此本段仅据材料确认事件的时序、行动和可见后果，并把政策文本与地方实践、报称数字和社会经历分开处理。

\eventanchor{ML-1521-002}{明·1521（正德十六年）四月21日}
\paragraph{1521（正德十六年）四月21日：托梅·皮雷的政党被驱逐出北京} 这项记录把本章所见的制度运作落实到具体年份。其形成背景是：继承、官僚协调或皇权运作的压力使问题进入朝廷程序。 它带来的、而且仍可由后续材料追踪的改变是：它影响后续人事与政策议程；动机和派系标签不能由单一叙事裁定。 《武宗实录》正德十六年条；《明史·本纪》及相关列传；诏令、奏疏或会典版本 提供相互检验的起点；但译写候选以编年材料定位；实录记载反映朝廷选择，崇祯期须另以奏疏、地方及敌对政权材料交叉。 因此本段仅据材料确认事件的时序、行动和可见后果，并把政策文本与地方实践、报称数字和社会经历分开处理。

\eventanchor{ML-1521-003}{明·1521（正德十六年）四月27日}
\paragraph{1521（正德十六年）四月27日：朱厚熜即位嘉靖皇帝} 这项记录把本章所见的制度运作落实到具体年份。其形成背景是：继承、官僚协调或皇权运作的压力使问题进入朝廷程序。 它带来的、而且仍可由后续材料追踪的改变是：它影响后续人事与政策议程；动机和派系标签不能由单一叙事裁定。 《武宗实录》正德十六年条；《明史·本纪》及相关列传；诏令、奏疏或会典版本 提供相互检验的起点；但译写候选以编年材料定位；实录记载反映朝廷选择，崇祯期须另以奏疏、地方及敌对政权材料交叉。 因此本段仅据材料确认事件的时序、行动和可见后果，并把政策文本与地方实践、报称数字和社会经历分开处理。

\eventanchor{ML-1521-004}{明·1521（正德十六年）五月}
\paragraph{1521（正德十六年）五月：屯门之战：明军将拒绝离开的葡萄牙舰队驱逐出屯门} 这项记录把本章所见的制度运作落实到具体年份。其形成背景是：朝贡名义、边境安全与港口贸易的利益使交涉持续发生。 它带来的、而且仍可由后续材料追踪的改变是：它为后续册封、互市或海防安排留下文书与跨政权比较线索。 《武宗实录》正德十六年条；《明史·外国传》；《朝鲜王朝实录》、琉球《历代宝案》或海外档案 提供相互检验的起点；但译写候选以编年材料定位；实录记载反映朝廷选择，崇祯期须另以奏疏、地方及敌对政权材料交叉。 因此本段仅据材料确认事件的时序、行动和可见后果，并把政策文本与地方实践、报称数字和社会经历分开处理。

\eventanchor{ML-1521-005}{明·1521（正德十六年）}
\paragraph{1521（正德十六年）：宫殿重建完成} 这项记录把本章所见的制度运作落实到具体年份。其形成背景是：继承、官僚协调或皇权运作的压力使问题进入朝廷程序。 它带来的、而且仍可由后续材料追踪的改变是：它影响后续人事与政策议程；动机和派系标签不能由单一叙事裁定。 《武宗实录》正德十六年条；《明史·本纪》及相关列传；诏令、奏疏或会典版本 提供相互检验的起点；但译写候选以编年材料定位；实录记载反映朝廷选择，崇祯期须另以奏疏、地方及敌对政权材料交叉。 因此本段仅据材料确认事件的时序、行动和可见后果，并把政策文本与地方实践、报称数字和社会经历分开处理。

\eventanchor{MM-1521-TUNMEN-CONFLICT}{明·1521（正德十六年）}
\paragraph{1521（正德十六年）：屯门海域发生明军与葡萄牙船队冲突} 这项记录把本章所见的制度运作落实到具体年份。其形成背景是：葡萄牙使团在华身份和贸易要求受满剌加问题影响，广东官府同时担忧未经许可的停泊和武装活动。 它带来的、而且仍可由后续材料追踪的改变是：冲突未终止葡萄牙人对华贸易的尝试，但使后续居留与交易须通过地方官府和中介反复协商。 《明史·佛郎机传》；《广东通志》相关条目；葡萄牙航海记载 提供相互检验的起点；但冲突发生及其大致时段可由中葡材料互证；舰船数、胜负叙述和具体海域在不同记录中不一。 因此本段仅据材料确认事件的时序、行动和可见后果，并把政策文本与地方实践、报称数字和社会经历分开处理。
